\section{Transitive Sets}

This is a quick recap section on transitive sets and ordinals. We come to cardinals in \Cref{Ch1:Sec:Cardinals}.

\begin{boxdefinition}[Transitive Set]
    A set $A$ is \textbf{transitive} if $\forall y \in A \brac{y \subseteq A \lr \forall y \in A \, \forall x \in y \, \parenth{x \in A}}$
\end{boxdefinition}

\subsection{Ordinals}

\begin{boxdefinition}[Ordinal]
    A set $\alpha$ is an \textbf{ordinal} iff $\alpha$ is a transitive set and $\parenth{\alpha, \in}$ is a well-ordering.
\end{boxdefinition}

Note that writing $\parenth{A, \in}$ really means the structure $(A, R)$ where $R = \setst{(x, y) \in A \times A}{x \in y}$. It's slightly problematic to call $\in$ itself a relation because it isn't a set: it's (strictly speaking) just a symbol, which is interpreted as a set (a relation) in a model of set theory. We have to be careful with these sorts of things.

\begin{boxexample}[Examples of Ordinals]\hfill
    \begin{enumerate}
        \item $0 = \emptyset$, $1 = \set{0}$, $2 = \set{0, 1}$, $3 = \set{0, 1, 2}$, etc.
        \item $\omega = \set{0, 1, 2, \ldots}$
        \item $\omega + 1 = \omega \cup \set{\omega}$, $\omega + 2$, etc.
        \item $\omega + \omega$
    \end{enumerate}
    and so on.
\end{boxexample}

\subsection{Mostowski Collapse}

We will eventually prove this in more generality, but here's a useful theorem.

\begin{boxtheorem}[Mostowski Collapse Theorem for Well-Orderings]
    For all well-orderings $(A, R)$, there is a unique ordinal $\alpha$ such that there is a unique isomorphism $\pi$ from $(A, R)$ to $(\alpha, \in)$ given by
    \begin{align*}
        \pi(y) = \setst{\pi(x)}{x R y}
    \end{align*}
\end{boxtheorem}

This is a special case of what's called (surprise surprise) the Mostowski Collapse Theorem. Some things we can do include removing the well/linear requirement on the ordering and replace ``ordinal'' with just ``transitive set''.

What the theorem above says is that the ordinals form a complete set of representatives for the class of all well-orderings up to isomorphism.
